\documentclass[12pt]{article}
\usepackage{fullpage}
\usepackage{amsmath,amssymb,amsthm}

\newtheorem{lemma}{Lemma}
\newtheorem{proposition}{Proposition}
\newcommand{\Dil}{\operatorname{Dil}}
\newcommand{\Vol}{\operatorname{Vol}}
\newcommand{\Mass}{\operatorname{Mass}}
\newcommand{\Fill}{\operatorname{Fill}}
\newcommand{\EFill}{\operatorname{EFill}}
\newcommand{\RR}{\mathbb R}

\begin{document}

\begin{center}
{\bf Rigorous partial proof and the remaining quotient tightening estimate}
\end{center}

\section*{Problem statement and interpretation}

The small universal constant in the problem is denoted by \(\kappa\), to avoid
confusion with \(k\)-dilation.  Rectangles are oriented Euclidean products, and
degree \(1\) means relative degree \(1\) on
\(H_4(R,\partial R;\mathbb Z)\to H_4(S,\partial S;\mathbb Z)\).

This document proves the first alternative and the strongest scalar hard-case
bound obtained from the standard estimates.  It also gives a precise quotient
form of the averaged tightening estimate which would finish the theorem.  That
quotient estimate is not proved here; thus the document is still a conditional
reduction, not a complete proof of the stated dichotomy.  The arguments are
written for smooth transverse maps.  For piecewise smooth maps one first
approximates relative to the boundary strata by smooth maps with
\(\Dil_2\le1+\varepsilon\) and the same relative degree, applies the estimates
to the target scaled by \((1+\varepsilon)^{-1/2}\), and finally decreases
\(\kappa\) by a universal factor and lets \(\varepsilon\downarrow0\).

\section*{Standard inputs}

All constants below are positive and universal.  If \(\Dil_2(f)\le1\), then
\(\Dil_3(f),\Dil_4(f)\le1\): for singular values
\(\sigma_1\ge\cdots\ge\sigma_4\), the inequality
\(\sigma_1\sigma_2\le1\) implies \(\sigma_2\le1\) and
\(\sigma_1\cdots\sigma_j\le\sigma_1\sigma_2\) for \(j\ge2\).

We use Guth's map-of-pairs estimate \cite[Estimate 1]{GuthExpanding}.  For
\(n=4,k=2,j=1,l=4\), applied to a degree-one map with \(\Dil_2\le1\), it gives
\begin{equation}\label{guth14}
        R_1^2\,\Vol(R)\ge c_0\,S_1^3S_2S_3S_4 .
\end{equation}
We also use Guth's small-volume relative isoperimetric estimate in rectangles
\cite[Theorem 3]{GuthExpanding}: if \(Z\) is a relative \(2\)-cycle in \(R\)
with \(A=\Mass Z\le c_*R_1R_2\), then
\[
        \Fill_R(Z)\le C_*\max\{A^{3/2},A^2/R_1\}.
\]
By the contrapositive, after decreasing constants,
\begin{equation}\label{profile-inverse}
\Fill_R(Z)\ge L\quad\Longrightarrow\quad
\Mass Z\ge c\min\{R_1R_2,L^{2/3},(R_1L)^{1/2}\}.
\end{equation}
In particular,
\begin{equation}\label{iso}
        \Mass Z<c_*R_1R_2\quad\Longrightarrow\quad
        \Fill_R(Z)\le C R_1R_2^2 .
\end{equation}

\section*{The first alternative}

Let \(Q\) be the concentric subrectangle of the \((x_3,x_4)\)-face of \(S\)
with side lengths \(S_3/2,S_4/2\), and put
\[
        P_y=[0,S_1]\times[0,S_2]\times\{y\},\qquad y\in Q.
\]

\begin{lemma}[central target plane]\label{central}
Every relative \(3\)-chain in \(S\) with boundary \(P_y\), \(y\in Q\), has
mass at least \(cS_1S_2S_3\).
\end{lemma}

\begin{proof}
Choose a \(1\)-Lipschitz function \(\psi\) on the
\((x_3,x_4)\)-rectangle which vanishes on its boundary and satisfies
\(\psi\ge cS_3\) on \(Q\).  Put \(\beta=\psi\,dx_1\wedge dx_2\).  Then
\(d\beta\) has comass at most \(1\), and Stokes' theorem gives
\[
\Mass Y\ge\left|\int_Yd\beta\right|
 =\left|\int_{P_y}\beta\right|\ge cS_1S_2S_3 .
\]
The relative boundary terms vanish because \(dx_1\wedge dx_2=0\) on the
\(x_1,x_2\)-faces and \(\psi=0\) on the \(x_3,x_4\)-faces.
\end{proof}

\begin{proposition}\label{firstalt}
There is \(c_2>0\) such that, under the hypotheses of the problem, if
\[
        R_1R_2^2\le c_2S_1S_2S_3 ,
\]
then \(R_3R_4\ge c_2S_3S_4\).
\end{proposition}

\begin{proof}
Let \(T=[[R]]\) be the fundamental relative current and set
\(F=(f_3,f_4)\).  Since \(f\) has relative degree \(1\), the constancy theorem
gives \(f_\#T=[[S]]\) as a relative \(4\)-current.  For a.e. \(y\in Q\),
slicing naturality gives
\[
        f_\#\langle T,F,y\rangle
        =\langle [[S]],(x_3,x_4),y\rangle=\pm P_y .
\]
Write \(Z_y=\langle T,F,y\rangle\).  Any relative filling of \(Z_y\) in \(R\)
pushes forward, using \(\Dil_3(f)\le1\), to a filling of \(P_y\) in \(S\).
Lemma \ref{central} therefore gives
\[
        \Fill_R(Z_y)\ge cS_1S_2S_3 .
\]
If \(c_2\) is sufficiently small, \eqref{iso} forces
\(\Mass Z_y\ge c_*R_1R_2\).  Since \(J_2F\le\Dil_2(f)\le1\), coarea gives
\[
 \Vol(R)\ge\int_Q\Mass Z_y\,dy
        \ge cR_1R_2S_3S_4 .
\]
Cancelling \(R_1R_2\) proves the proposition.
\end{proof}

\section*{The unconditional scalar hard-case bound}

Assume \(R_1=\alpha S_1\le S_1\) and that the first conclusion fails with a
small constant.  Proposition \ref{firstalt} gives
\begin{equation}\label{r12large}
        R_1R_2^2\ge cS_1S_2S_3 .
\end{equation}
For the slices above, \eqref{profile-inverse} with
\(L=cS_1S_2S_3\) gives
\[
\Mass Z_y\ge c\min\{R_1R_2,(S_1S_2S_3)^{2/3},
                         (R_1S_1S_2S_3)^{1/2}\}.
\]
Using \eqref{r12large} and \(R_1=\alpha S_1\), each term in the minimum is
at least \(c\alpha^{1/2}S_1(S_2S_3)^{1/2}\).  Hence
\begin{equation}\label{slicevol}
        \Vol(R)\ge
        c\alpha^{1/2}S_1S_2^{1/2}S_3^{3/2}S_4 .
\end{equation}
On the other hand, \eqref{guth14} gives
\begin{equation}\label{monovol}
        \Vol(R)\ge
        c\alpha^{-2}S_1S_2S_3S_4 .
\end{equation}
Optimizing between \eqref{slicevol} and \eqref{monovol} yields
\[
        \Vol(R)\ge cS_1S_2^{3/5}S_3^{7/5}S_4 .
\]
Thus purely scalar slicing and monomial estimates miss the desired hard-case
scale by the factor \((S_2/S_3)^{1/10}\).

\section*{A quotient averaged estimate sufficient for the theorem}

Keep \(F=(f_3,f_4)\), put \(G=(f_1,f_2)\), and define
\[
E_y=\int_{Z_y}\|dG|_{TZ_y}\|_{HS}^2\,d\|Z_y\|,\qquad
I=\int_QE_y\,dy,\qquad q=|Q|.
\]
At a regular point of \(F\), let \(K=\ker dF\), and let \(a\ge b\) be the
singular values of \(dF|_{K^\perp}\).  For a unit \(w\in K\) and a unit
\(v\in K^\perp\) with \(|dF(v)|=a\), the mixed \(G\)-\(F\) component of
\(df(w)\wedge df(v)\) has norm \(|dG(w)|a\).  Therefore
\(|dG(w)|a\le1\), and
\[
        \|dG|_{\ker dF}\|_{HS}^2J_2F\le2 .
\]
Coarea gives
\begin{equation}\label{Icoarea}
        I\le C\Vol(R).
\end{equation}
Let \(\Omega=\{y\in Q:E_y\le4I/q\}\); then \(|\Omega|\ge3q/4\).

To avoid the null-chain obstruction, work in the quotient complex
\[
 {\mathcal Q}_*=N_*(R,\partial R;\mathbb R)/\ker f_\# .
\]
For a central regular \(y\), define the relaxed essential filling norm
\[
\EFill_f(y)=\inf\{\Mass C:\ C\in N_3(R,\partial R;\mathbb R),\
        f_\#(\partial C)=\pm P_y\}.
\]
This is the filling norm of the degree-carrying class in \({\mathcal Q}_*\).
If \(C\) is admissible, then \(f_\#C\) fills \(P_y\) in \(S\), so Lemma
\ref{central} and \(\Dil_3(f)\le1\) imply
\begin{equation}\label{AElower}
A_\Omega:=\int_\Omega\EFill_f(y)\,dy
        \ge cS_1S_2S_3\,q .
\end{equation}

For \(z\in Q\), let \({\mathcal B}_z\) be the set of points \(x\in R\) such
that \(z\) belongs to the \(F\)-image of at least one coordinate \(2\)-plaque
through \(x\).  For fixed \(x\), the set of such \(z\)'s has measure at most
\(CR_3R_4\), because there are six plaques and each has \(F\)-area at most
\(R_3R_4\).  Consequently, for every normal \(3\)-current \(B\),
\begin{equation}\label{badfubini}
 \frac1q\int_Q\Mass(B\llcorner{\mathcal B}_z)\,dz
 \le C\frac{R_3R_4}{S_3S_4}\Mass B .
\end{equation}
For a quotient class \([C]\in{\mathcal Q}_3\), set
\[
 \beta_z([C])=\inf_{N\in\ker f_\#}\Mass((C+N)\llcorner{\mathcal B}_z).
\]
Then \(\beta_z([C])\le \Mass(C\llcorner{\mathcal B}_z)\).

The following quotient averaged tightening estimate would finish the proof:
for measurable near-minimizers \(C_y^0\), \(\Mass C_y^0\le2\EFill_f(y)\),
\begin{equation}\tag{QEAT}\label{QEAT}
 A_\Omega
 \le C\left(R_1I+\frac{I^2}{S_1q}\right)
 +\frac{C}{q}\int_Q\int_\Omega
       \beta_z([C_y^0])\,dy\,dz .
\end{equation}
Indeed, by \(\beta_z([C_y^0])\le\Mass(C_y^0\llcorner{\mathcal B}_z)\) and
\eqref{badfubini},
\[
A_\Omega\le C\left(R_1I+\frac{I^2}{S_1q}\right)
  +C\frac{R_3R_4}{S_3S_4}A_\Omega .
\]
If \(R_3R_4\le\kappa S_3S_4\), the last term is absorbed for small
\(\kappa\).  If also \(R_1\le\kappa S_1\) and
\(I\le\delta S_1S_2^{1/2}S_3^{3/2}S_4\), then, since \(q\sim S_3S_4\),
\[
R_1I\le C\kappa\delta S_1S_2S_3^2S_4,\qquad
\frac{I^2}{S_1q}\le C\delta^2S_1S_2S_3^2S_4 ,
\]
contradicting \eqref{AElower} for small \(\delta,\kappa\).  Thus
\(I\ge cS_1S_2^{1/2}S_3^{3/2}S_4\), and \eqref{Icoarea} gives the desired
second alternative.

\section*{Remaining open issues}

The document is still a conditional reduction.

\begin{enumerate}
\item The essential missing estimate is \eqref{QEAT}.  It is a quotient or dual
averaged tightening theorem for the degree-carrying part of the slices
\(Z_y=\langle[[R]],(f_3,f_4),y\rangle\).

\item The quotient formulation is needed.  In a two-cell model with
\(f_\#A=f_\#A'=P\), the difference \(A-A'\) is \(f\)-null, but a source-local
restriction to a bad set containing \(A\) and not \(A'\) is no longer
\(f\)-null.  Thus an ordinary source-local deformation does not descend to
\({\mathcal Q}_*\).

\item The proved parts are the first alternative, the scalar partial hard-case
bound, the weighted coarea inequality, the essential filling lower bound, the
exact-plaque Fubini estimate, and the absorption algebra following
\eqref{QEAT}.

\item A complete proof would require a quotient-equivariant exact-plaque
tightening, or equivalently a dual target-pullback cochain homotopy, proving
\eqref{QEAT}.
\end{enumerate}

\begin{thebibliography}{9}

\bibitem{GuthExpanding}
L. Guth, \emph{Area-expanding embeddings of rectangles}, arXiv:0710.0403.

\bibitem{GuthThesis}
L. Guth, \emph{Area-contracting maps between rectangles}, Ph.D. thesis,
Massachusetts Institute of Technology, 2005.

\bibitem{GuthDirectional}
L. Guth, \emph{Directional isoperimetric inequalities and rational homotopy
invariants}, arXiv:0802.3549.

\end{thebibliography}

\end{document}
